\documentclass[11pt]{article}
\usepackage[utf8]{inputenc}
\usepackage{amsmath,amssymb}
\usepackage[margin=1in]{geometry}
\usepackage{hyperref}
\emergencystretch=3em

\title{Joint continuity of the amplitude coefficients and the measurable chart integral}
\author{Claude, with Daniel Murfet}
\date{2026-09-07}

\begin{document}
\maketitle
\sloppy

\begin{abstract}
Infrastructure for the probabilistic normalised-remainder theorem (Astra~\#11 rank~4, second half).
The amplitude coefficients $c_k(s)=\texttt{ampCoeff}\,\beta\,x\,y\,k\,s$ are finite algebraic
expressions in finitely many entries of $(x,y)$ and in $s$, hence jointly continuous on
$\texttt{CoeffPair}\times\mathbb R$ (\texttt{ampCoeff\_pair\_continuous}). For a measurable random
Taylor datum $X\colon\Omega\to\texttt{CoeffPair}$ the chart integral $Z_N(X(\omega))$ is therefore a
measurable function of $\omega$ (\texttt{measurable\_chartZ\_comp}), by the parameter-measurability
theorem of unit~138. The chart integral is packaged as \texttt{chartZ} on \texttt{CoeffPair}, and the
uniform Taylor tree of unit~135 is restated on norm balls with the coefficient maps
\texttt{coeffA}, \texttt{coeffB} (\texttt{chartZ\_taylor\_tree\_ball}). Zero \texttt{sorry}/\texttt{axiom}.
\end{abstract}

\section{The things formalised}
{\small
\begin{verbatim}
theorem continuous_eval_coeffPair (i : PairIdx) : Continuous fun a : CoeffPair => a i
theorem continuous_toX / continuous_toY (ρ) (k) :
    Continuous fun a : CoeffPair => toX ρ a k
theorem convPow_continuous_param (f : X → ℕ × ℕ → ℝ)
    (hf : ∀ k, Continuous fun p => f p k) (n) :
    ∀ k, Continuous fun p => convPow (f p) n k
theorem dropConst_continuous_param / expCoeff_continuous_param
theorem ampCoeff_continuous_param (β) (x y : X → ℕ × ℕ → ℝ) (hx hy) (s : X → ℝ) (hs) (k) :
    Continuous fun p => ampCoeff β (x p) (y p) k (s p)
theorem ampCoeff_pair_continuous (β ρ) (k) :
    Continuous fun q : CoeffPair × ℝ => ampCoeff β (toX ρ q.1) (toY ρ q.1) k q.2
theorem ampCoeff_pair_measurable (X : Ω → CoeffPair) (hX : Measurable X) (k) :
    Measurable fun p : Ω × ℝ => ampCoeff β (toX ρ (X p.1)) (toY ρ (X p.1)) k p.2
def chartZ (β b ρ) (h₁ h₂ k₁ k₂) (N) (a : CoeffPair) : ℝ :=
  twoDAmp β b N h₁ h₂ k₁ k₂ (anaAmp (ampCoeff β (toX ρ a) (toY ρ a)) b)
theorem measurable_chartZ_comp ... (X) (hX : Measurable X) :
    Measurable fun ω => chartZ … N (X ω)
theorem measurable_chartZ ... : Measurable (chartZ β b ρ h₁ h₂ k₁ k₂ N)
theorem chartZ_taylor_tree_ball ... (M) (hM : 0 ≤ M) (a) (ha : ‖a‖ ≤ M) (T N)
    (hN : 1 ≤ N) :
    |chartZ … N a - ∑ α ∈ polesBelowGen h₁ h₂ k₁ k₂ p₁ p₂ T,
        N^(-α) * (coeffA β ρ … α a * log N + coeffB β b ρ … α a)|
      ≤ uniformTreeConst β b ρ M (2 * M) p₁ p₂ T h₁ h₂ k₁ k₂ 0 * (N^(-(2 * T)) * (1 + log N))
\end{verbatim}
}

\section{Mathematics}
Each $c_k(s)$ is $e^{\beta sx_{00}}\sum_{a\le k}y_a\,E_{k-a}(\beta s)$ with
$E_m(t)=\sum_{n\le|m|}\frac{t^n}{n!}(\bar x^{*n})_m$, and $(\bar x^{*n})_m$ is a finite sum of
products of entries of $\bar x$; so $c_k$ is a polynomial in finitely many entries of $(x,y)$ and in
$e^{\beta sx_{00}}$, $s$. Point evaluation is $1$-Lipschitz on $\ell^1$, so all of this is jointly
continuous on $\texttt{CoeffPair}\times\mathbb R$, hence Borel measurable (the product of a metric
space with the second-countable $\mathbb R$ has Borel $\sigma$-algebra equal to the product
$\sigma$-algebra). Composing with a measurable $X$ gives the joint measurability hypothesis of
unit~138's parameter theorem, whose conclusion is the measurability of $\omega\mapsto Z_N(X(\omega))$.
On the ball $\|a\|\le M$ the amplitude has the common envelope $Me^{2\beta Ms}$ (unit~145), so the
uniform Taylor tree applies with $C_0=M$, $L=2M$, $D=0$.

\section{Paper-facing meaning}
Two of the three ingredients of the probabilistic normalised-remainder statement are now in place:
the deterministic bounds (unit~148, applied with $K=\texttt{uniformTreeConst}$ from
\texttt{chartZ\_taylor\_tree\_ball}) and the measurability of the random chart integral. The next
unit assembles them with Slutsky's theorem: for random Taylor data bounded by $M$ and converging in
distribution, $(Z_n-L_{<\alpha})/(N_n^{-\alpha}\log N_n)\Rightarrow A_\alpha(Z)$ and the B-slot
analogue.

\section{Lean notes}
\begin{itemize}
\item Continuity by induction on the convolution power: the \texttt{zero} case is
\texttt{change Continuous fun \_ => delta k} (definitional unfolding of \texttt{convPow a 0}),
the \texttt{succ} case unfolds \texttt{conv} to a \texttt{continuous\_finsetSum}.
\item Binders over a product must be type-annotated (\texttt{fun q : CoeffPair × ℝ => …});
otherwise projections \texttt{q.1} fail to elaborate (``Invalid projection'').
\item \texttt{lp.lipschitzWith\_one\_eval 1 i} gives continuity of point evaluation on $\ell^1$;
\texttt{Continuous.measurable} on \texttt{CoeffPair × ℝ} finds the product
\texttt{OpensMeasurableSpace} instance through \texttt{SecondCountableTopologyEither} with the
$\mathbb R$ factor.
\item The file compiled after one binder-annotation fix.
\end{itemize}

\end{document}
